\documentclass[conference]{IEEEtran}
\usepackage[utf8]{inputenc}
\usepackage[T1]{fontenc}
\usepackage{amsmath}
\usepackage{amsfonts}
\usepackage{amssymb}
\usepackage{graphicx}
\usepackage{hyperref}
\usepackage{listings}
\usepackage{xcolor}
\usepackage{url}
\usepackage{cite}

% Hyperref settings
\hypersetup{
    colorlinks=true,
    linkcolor=blue,
    filecolor=magenta,
    urlcolor=cyan,
    citecolor=green,
    pdftitle={Design and Analysis of a Multimodal AI Chatbot with Intelligent Resume Building and ATS Optimization},
    pdfauthor={Research Paper},
    pdfsubject={AI Chatbot, Multi-Modal Processing, NLP},
    pdfkeywords={AI Chatbot, Multi-Modal Processing, Natural Language Processing, Web Application, Resume Builder, Voice Recognition}
}

% Code listing style
\lstset{
    basicstyle=\ttfamily\small,
    breaklines=true,
    frame=single,
    numbers=left,
    numberstyle=\tiny\color{gray},
    keywordstyle=\color{blue},
    commentstyle=\color{green},
    stringstyle=\color{red}
}

% Title
\title{Design and Analysis of a Multimodal AI Chatbot with Intelligent Resume Building and ATS Optimization}

\author{\IEEEauthorblockN{Research Paper}
\IEEEauthorblockA{AI Chatbot Project\\Version 1.0, 2024}}

\begin{document}

\maketitle

\begin{abstract}
This paper presents the design and analysis of a multimodal, web-based AI chatbot system that integrates conversational assistance with intelligent resume building and Applicant Tracking System (ATS) optimization. The system leverages Google's Gemini 2.0 Flash API for natural language processing, incorporates multimodal input support including images, videos, PDF documents, and text files, and features an integrated resume builder with AI-powered content generation. Methodologically, the work adopts an architecture-centric review and system analysis approach, examining component interactions, data flows, and technology choices. The study qualitatively evaluates system functionality, usability, and expected scalability characteristics by comparing implemented capabilities with those reported in existing chatbot and recruitment automation systems. The primary contribution is a unified system architecture and design framework that demonstrates how conversational AI, multimodal processing, and ATS-aware resume generation can be integrated within a single web-based platform for career development applications.
\end{abstract}
\begin{IEEEkeywords}
AI Chatbot, Multi-Modal Processing, Natural Language Processing, Web Application, Resume Builder, Voice Recognition
\end{IEEEkeywords}

\section{Introduction}

\subsection{Background}

Artificial Intelligence (AI) chatbots have revolutionized human-computer interaction by providing intelligent, context-aware responses to user queries. Modern chatbot systems have evolved from simple rule-based systems to sophisticated neural network-based models capable of understanding context, processing multiple data types, and maintaining conversational history \cite{adamopoulou2020}. The integration of multi-modal capabilities allows chatbots to process not only text but also images, videos, and documents, significantly expanding their utility \cite{radford2021}.

\subsection{Problem Statement}

Traditional chatbot systems often lack comprehensive file processing capabilities, limited export functionalities, and insufficient mechanisms for preserving conversation history. Additionally, the integration of specialized tools such as resume builders within chatbot interfaces remains underexplored. Existing studies rarely examine the integration of conversational AI with ATS-optimized resume generation within a single web-based system, leading to fragmented workflows and duplicated design efforts. This project addresses these limitations by developing a unified platform that combines conversational AI with advanced file processing and ATS-aware document generation capabilities.

\subsection{Objectives}

The primary objectives of this review are to:
\begin{enumerate}
    \item Analyze the existing architecture and implementation of the AI chatbot system
    \item Evaluate current features and their effectiveness
    \item Identify areas for improvement and propose new features
    \item Provide a comprehensive reference for researchers and developers
    \item Assess system performance, usability, and scalability limitations from a research perspective
\end{enumerate}

\section{Literature Review}

\subsection{Evolution of Chatbot Technology}

Chatbot technology has undergone significant evolution since the development of ELIZA in the 1960s \cite{weizenbaum1966}. Modern chatbots utilize deep learning architectures, particularly transformer-based models, which have shown remarkable performance in natural language understanding and generation tasks \cite{vaswani2017}. The introduction of large language models (LLMs) such as GPT series and Gemini has further enhanced chatbot capabilities \cite{thoppilan2022}. However, much of this work focuses on generic conversational ability and open-domain dialog, leaving limited guidance for systems tightly integrated into specialized workflows such as resume creation. The present system extends this line of research by embedding LLM-based conversational capabilities within a resume-oriented interaction flow and analyzing system-level trade-offs.

\subsection{Multi-Modal AI Systems}

Multi-modal AI systems process and integrate information from multiple sources including text, images, audio, and video \cite{baltrusaitis2019}. Research has shown that combining visual and textual information significantly improves AI system performance in understanding context and generating appropriate responses \cite{lu2019}. The integration of PDF processing and document analysis in chatbot systems has been explored in various studies \cite{devlin2018}. Nevertheless, existing multimodal systems are often evaluated on benchmark tasks rather than on career-focused use cases that combine resumes, job descriptions, and supporting documents. This study targets multimodal processing specifically for user-provided career materials, linking document understanding directly to conversational guidance and resume generation.

\subsection{Voice Recognition in Web Applications}

Web-based voice recognition has become increasingly accessible through browser APIs such as the Web Speech API \cite{w3c2023}. Studies have demonstrated the effectiveness of voice input in improving user experience and accessibility in web applications \cite{pradhan2018}. The integration of speech recognition with chatbot systems enables more natural human-computer interaction, yet prior work often treats speech input as an alternative modality for simple query entry. This leaves limited analysis of voice interaction within complex, stateful workflows that involve document uploads, resume editing, and ATS optimization. The proposed system contributes to this space by combining browser-based voice recognition with persistent chat history and document-centric interactions.

\subsection{Resume Generation and ATS Optimization}

Automated resume generation systems have gained prominence with the increasing use of Applicant Tracking Systems (ATS) \cite{black2020}. Recent work on AI-driven resume analysis and ATS feedback loops suggests growing academic attention to resume parsing and optimization \cite{resumai2024,atsllm2024,resumebuilder2024}. However, many existing systems are implemented as standalone resume builders or ATS scoring tools, offering limited conversational support and weak integration with broader career assistance workflows. This study addresses that limitation by embedding ATS-aware resume generation inside a multimodal chatbot, enabling iterative, dialog-based refinement rather than one-off, form-based generation.

\section{System Architecture and Existing Features}

\subsection{System Overview}

The system is implemented as a client-side web application using HTML5, CSS3, and JavaScript, with integration of external APIs for AI processing. The architecture follows a modular design pattern, separating concerns between user interface, data processing, and API communication.

\subsection{Core Technologies}

\begin{itemize}
    \item \textbf{Frontend Framework:} Vanilla JavaScript with modern ES6+ features
    \item \textbf{AI API:} Google Gemini 2.0 Flash (Generative Language API) \cite{google2024}
    \item \textbf{PDF Processing:} PDF.js library for client-side PDF text extraction \cite{mozilla2024}
    \item \textbf{Voice Recognition:} Web Speech API (WebKit Speech Recognition)
    \item \textbf{Storage:} LocalStorage API for client-side data persistence
    \item \textbf{Styling:} Custom CSS with CSS Variables for theme management
\end{itemize}
From a research and design perspective, a predominantly client-side architecture was selected to simplify deployment, reduce backend complexity, and keep user data within the user's browser, at the cost of limited scalability and constrained data persistence. Google Gemini 2.0 Flash was preferred over alternative LLMs because of its native multimodal support, competitive latency for interactive applications, and ease of integration through a single generative API. LocalStorage was chosen instead of a server-side database to enable a fully static, easily hostable prototype that preserves privacy, while explicitly accepting browser storage limits and the lack of centralized analytics as key trade-offs.

\subsection{Existing Features}

\subsubsection{Conversational AI Interface}
\begin{itemize}
    \item Real-time text-based conversation with AI
    \item Typing effect animation for natural response display
    \item Context-aware conversation history maintenance
    \item Response streaming and interruption capabilities
\end{itemize}

\subsubsection{Multi-Modal File Processing}
\begin{itemize}
    \item \textbf{Image Processing:} Support for various image formats (JPEG, PNG, GIF, WebP)
    \item \textbf{Video Processing:} MP4 and WebM video file support
    \item \textbf{PDF Processing:} Automatic text extraction from PDF documents using PDF.js
    \item \textbf{Text File Processing:} Support for .txt and .csv files with content extraction
    \item \textbf{File Preview:} Real-time preview of uploaded files before processing
\end{itemize}

\subsubsection{Chat Management Features}
\begin{itemize}
    \item \textbf{Auto-Save:} Automatic saving of chat history to browser localStorage
    \item \textbf{Manual Save:} Export chat history as JSON files
    \item \textbf{Load Functionality:} Import and restore previously saved chat sessions
    \item \textbf{Export Options:} Export chat history as text files or PDF documents
    \item \textbf{Search Capability:} Full-text search within chat history with highlighting
    \item \textbf{Copy Functionality:} Double-click to copy individual messages
\end{itemize}

\subsubsection{Voice Input Integration}
\begin{itemize}
    \item Browser-based speech recognition
    \item Real-time transcription to text input
    \item Visual feedback during recording (pulsing animation)
    \item Support for multiple languages (configurable)
\end{itemize}

\subsubsection{Resume Builder Module}
\begin{itemize}
    \item \textbf{Template Selection:} Multiple resume templates (Clean ATS, Bold Header, Modern Blue, Creative Layout)
    \item \textbf{AI-Powered Content Generation:}
    \begin{itemize}
        \item Professional summary generation based on job title
        \item Skills recommendation using Cohere API \cite{cohere2024}
        \item ATS optimization suggestions
    \end{itemize}
    \item \textbf{Comprehensive Sections:}
    \begin{itemize}
        \item Personal information
        \item Professional summary
        \item Education details
        \item Work experience
        \item Skills and certifications
        \item Projects portfolio
        \item Languages and interests
        \item Achievements and volunteer experience
    \end{itemize}
    \item \textbf{Export Functionality:} Print or save resume as PDF/Word document
\end{itemize}
Beyond its practical functionality, the resume builder module represents a design contribution that demonstrates an applied AI use case within the broader chatbot architecture. In particular, the ATS optimization suggestions align resume structure and keyword usage with common recruitment screening practices, theoretically increasing the probability of shortlisting compared to non-optimized resumes.

\subsubsection{User Interface Features}
\begin{itemize}
    \item \textbf{Theme Toggle:} Light and dark mode support
    \item \textbf{Responsive Design:} Mobile-friendly interface
    \item \textbf{Suggestion Cards:} Pre-defined query suggestions for quick interaction
    \item \textbf{File Upload Interface:} Drag-and-drop and click-to-upload support
    \item \textbf{Visual Feedback:} Loading states, error messages, and success notifications
\end{itemize}

\subsubsection{Advanced Functionalities}
\begin{itemize}
    \item \textbf{Response Control:} Stop generation mid-response
    \item \textbf{Chat History Management:} Clear all conversations
    \item \textbf{Context Preservation:} Maintains conversation context across sessions
    \item \textbf{Error Handling:} Comprehensive error handling with user-friendly messages
\end{itemize}

\section{Proposed New Features and Enhancements}

From a research planning standpoint, the enhancements are organized into short-term extensions aligned with the current scope and long-term directions that broaden the application domain. This grouping clarifies which items are feasible for immediate research validation and which require broader infrastructural or organizational deployments.

\subsection{Short-Term (Scope-Aligned) Enhancements}

\subsubsection{Enhanced AI Capabilities}

\paragraph{Multi-Language Support}
\begin{itemize}
    \item \textbf{Implementation:} Integration of translation APIs (Google Translate API)
    \item \textbf{Benefit:} Enable conversations in multiple languages
    \item \textbf{Use Case:} International users and multilingual content processing
\end{itemize}

\paragraph{Sentiment Analysis}
\begin{itemize}
    \item \textbf{Implementation:} Real-time sentiment detection in user messages
    \item \textbf{Benefit:} Adaptive responses based on user emotional state
    \item \textbf{Use Case:} Career coaching, stress-aware guidance, and user support
\end{itemize}

\paragraph{Context-Aware Suggestions}
\begin{itemize}
    \item \textbf{Implementation:} Machine learning-based query suggestions
    \item \textbf{Benefit:} Personalized recommendations based on conversation history
    \item \textbf{Use Case:} Improved user experience and engagement
\end{itemize}

\subsubsection{Advanced File Processing}

\paragraph{OCR (Optical Character Recognition)}
\begin{itemize}
    \item \textbf{Implementation:} Integration of Tesseract.js or cloud OCR services
    \item \textbf{Benefit:} Extract text from images and scanned documents
    \item \textbf{Use Case:} Document digitization and accessibility
\end{itemize}

\paragraph{Audio File Processing}
\begin{itemize}
    \item \textbf{Implementation:} Speech-to-text conversion for audio files
    \item \textbf{Benefit:} Process voice recordings and audio content
    \item \textbf{Use Case:} Meeting notes, podcast transcription
\end{itemize}

\paragraph{Excel/Spreadsheet Processing}
\begin{itemize}
    \item \textbf{Implementation:} SheetJS library for spreadsheet parsing
    \item \textbf{Benefit:} Analyze and process Excel files
    \item \textbf{Use Case:} Data analysis and reporting
\end{itemize}

\paragraph{Code File Analysis}
\begin{itemize}
    \item \textbf{Implementation:} Syntax highlighting and code analysis
    \item \textbf{Benefit:} Understand and explain code files
    \item \textbf{Use Case:} Programming assistance and code review
\end{itemize}

\subsubsection{Enhanced Resume Builder}

\paragraph{Industry-Specific Templates}
\begin{itemize}
    \item \textbf{Implementation:} Templates tailored for specific industries
    \item \textbf{Benefit:} Better alignment with industry standards
    \item \textbf{Use Case:} Sector-specific job applications
\end{itemize}

\paragraph{Real-Time ATS Score}
\begin{itemize}
    \item \textbf{Implementation:} Live ATS compatibility scoring
    \item \textbf{Benefit:} Immediate feedback on resume optimization
    \item \textbf{Use Case:} Resume improvement and optimization
\end{itemize}

\paragraph{Cover Letter Generator}
\begin{itemize}
    \item \textbf{Implementation:} AI-powered cover letter generation
    \item \textbf{Benefit:} Complete application package creation
    \item \textbf{Use Case:} Comprehensive job application support
\end{itemize}

\subsection{Long-Term (Future Research) Directions}

\subsubsection{Collaboration Features}

\paragraph{Shared Chat Sessions}
\begin{itemize}
    \item \textbf{Implementation:} WebSocket-based real-time collaboration
    \item \textbf{Benefit:} Multiple users can participate in the same conversation
    \item \textbf{Use Case:} Team discussions and collaborative problem-solving
\end{itemize}

\paragraph{Chat Sharing}
\begin{itemize}
    \item \textbf{Implementation:} Generate shareable links for conversations
    \item \textbf{Benefit:} Easy sharing of conversations with others
    \item \textbf{Use Case:} Knowledge sharing and documentation
\end{itemize}

\subsubsection{Analytics and Insights}

\paragraph{Usage Analytics Dashboard}
\begin{itemize}
    \item \textbf{Implementation:} Track conversation patterns and usage statistics
    \item \textbf{Benefit:} User behavior insights and system optimization
    \item \textbf{Use Case:} Performance monitoring and improvement
\end{itemize}

\paragraph{Conversation Summarization}
\begin{itemize}
    \item \textbf{Implementation:} Automatic generation of conversation summaries
    \item \textbf{Benefit:} Quick overview of lengthy conversations
    \item \textbf{Use Case:} Documentation and review
\end{itemize}

\subsubsection{Security and Privacy}

\paragraph{End-to-End Encryption}
\begin{itemize}
    \item \textbf{Implementation:} Client-side encryption for sensitive data
    \item \textbf{Benefit:} Enhanced privacy and security
    \item \textbf{Use Case:} Handling confidential information
\end{itemize}

\paragraph{Data Anonymization}
\begin{itemize}
    \item \textbf{Implementation:} Automatic removal of personal information
    \item \textbf{Benefit:} Privacy protection in exported data
    \item \textbf{Use Case:} GDPR compliance and data protection
\end{itemize}

\subsubsection{Integration Capabilities}

\paragraph{API Integration}
\begin{itemize}
    \item \textbf{Implementation:} RESTful API for third-party integrations
    \item \textbf{Benefit:} Extensibility and integration with other systems
    \item \textbf{Use Case:} Enterprise applications and workflows
\end{itemize}

\paragraph{Browser Extension}
\begin{itemize}
    \item \textbf{Implementation:} Chrome/Firefox extension
    \item \textbf{Benefit:} Access chatbot from any webpage
    \item \textbf{Use Case:} Universal accessibility
\end{itemize}

\paragraph{Mobile Application}
\begin{itemize}
    \item \textbf{Implementation:} React Native or Flutter mobile app
    \item \textbf{Benefit:} Native mobile experience
    \item \textbf{Use Case:} On-the-go access and notifications
\end{itemize}

\section{Methodology}

\subsection{Development Approach}

The system was developed using an iterative development methodology, focusing on modular design and progressive enhancement. The implementation follows best practices for web development including:

\begin{itemize}
    \item \textbf{Separation of Concerns:} Clear separation between HTML structure, CSS styling, and JavaScript logic
    \item \textbf{Event-Driven Architecture:} Asynchronous event handling for user interactions
    \item \textbf{API Integration:} RESTful API communication with external services
    \item \textbf{Error Handling:} Comprehensive try-catch blocks and user feedback mechanisms
\end{itemize}

\subsection{Data Flow}

\begin{enumerate}
    \item \textbf{User Input Processing:}
    \begin{itemize}
        \item Text input or voice recognition
        \item File upload and processing
        \item Data validation and sanitization
    \end{itemize}
    
    \item \textbf{API Communication:}
    \begin{itemize}
        \item Request formatting according to API specifications
        \item Error handling and retry mechanisms
        \item Response parsing and validation
    \end{itemize}
    
    \item \textbf{Response Display:}
    \begin{itemize}
        \item Typing effect animation
        \item Message rendering
        \item Auto-scroll to latest message
    \end{itemize}
    
    \item \textbf{Data Persistence:}
    \begin{itemize}
        \item LocalStorage for temporary storage
        \item File export for permanent storage
        \item Chat history management
    \end{itemize}
\end{enumerate}

\subsection{Testing Strategy}

\begin{itemize}
    \item \textbf{Unit Testing:} Individual function testing
    \item \textbf{Integration Testing:} API and component integration
    \item \textbf{User Acceptance Testing:} Real-world usage scenarios
    \item \textbf{Cross-Browser Testing:} Compatibility across different browsers
\end{itemize}

\subsection{Research Validation Perspective}

While the methodology primarily follows software development best practices, it also supports research validation by enabling systematic analysis of architectural choices, multimodal processing flows, and user interaction patterns. An architecture-centric evaluation is appropriate for early-stage AI systems where constraints such as API quotas, deployment complexity, and privacy considerations limit large-scale user experiments. In this study, validation focuses on functional completeness, qualitative usability observations, and comparative reasoning against documented capabilities of existing chatbot and recruitment automation solutions rather than on benchmark datasets or controlled user trials. The absence of real-world recruitment data and standardized ATS benchmark scores is acknowledged as a limitation and is identified as an important direction for future empirical work.

\section{Implementation Details}

\subsection{Key Algorithms}

\subsubsection{Typing Effect Algorithm}
The system implements a word-by-word typing simulation algorithm with an interval of 40ms per word for optimal user experience. This creates a natural reading experience similar to human typing.

\subsubsection{File Processing Pipeline}
\begin{enumerate}
    \item File type detection
    \item Format-specific processing (PDF, image, video, text)
    \item Content extraction
    \item Base64 encoding for API transmission
\end{enumerate}

\subsubsection{Chat History Management}
\begin{itemize}
    \item Array-based storage for conversation context
    \item Automatic cleanup of temporary messages
    \item Efficient search algorithm with highlighting
\end{itemize}

\subsection{Performance Optimizations}

\begin{itemize}
    \item \textbf{Lazy Loading:} Load resources on demand
    \item \textbf{Debouncing:} Reduce API calls for search functionality
    \item \textbf{Caching:} Store frequently accessed data
    \item \textbf{Compression:} Optimize file sizes for faster loading
\end{itemize}

\section{Results and Discussion}

\subsection{Functional Evaluation}

The system successfully implements all core features with high reliability:
\begin{itemize}
    \item \textbf{File Processing:} Successfully handles multiple file formats
    \item \textbf{AI Integration:} Reliable API communication with error handling
    \item \textbf{User Interface:} Intuitive and responsive design
    \item \textbf{Data Management:} Efficient storage and retrieval mechanisms
\end{itemize}

\subsection{Performance Metrics}

\begin{itemize}
    \item \textbf{Response Time:} Average API response time: 2-5 seconds
    \item \textbf{File Processing:} PDF processing time: 1-3 seconds for standard documents
    \item \textbf{Storage Efficiency:} LocalStorage usage optimized for browser limitations
    \item \textbf{User Experience:} Smooth animations and responsive interactions
\end{itemize}
From a comparative perspective, these qualitative performance observations are consistent with reported response times for web-based LLM applications and client-side PDF processing pipelines. Although no direct baseline measurements against alternative chatbot or resume builder systems were conducted, the architecture is expected to offer lower latency for file handling than purely server-side solutions while incurring similar end-to-end delays for AI-generated responses due to external API constraints. A more rigorous quantitative comparison with baseline systems and controlled user studies is left as future work.

\subsection{Limitations and Challenges}

\begin{enumerate}
    \item \textbf{Browser Compatibility:} Some features require modern browsers
    \item \textbf{API Rate Limits:} External API limitations may affect usage
    \item \textbf{File Size Restrictions:} Large files may cause performance issues
    \item \textbf{Offline Functionality:} Limited offline capabilities
    \item \textbf{Lack of Large-Scale User Studies:} No controlled user study or A/B testing was conducted to quantitatively assess usability, task success rates, or user satisfaction.
    \item \textbf{Absence of Quantitative ATS Metrics:} The system was not evaluated against real recruitment pipelines or standardized ATS benchmarks, so claims about improved shortlisting probability remain theoretical.
\end{enumerate}
Due to API constraints, privacy considerations, and deployment limitations, large-scale user evaluation and integration with production ATS platforms were not feasible within the scope of this work. Consequently, the results focus on architectural behavior, feature completeness, and expected advantages relative to documented systems rather than on statistically validated performance improvements.

\subsection{Future Work}

\begin{enumerate}
    \item Implementation of proposed new features
    \item Performance optimization and scalability improvements
    \item Enhanced security measures
    \item Mobile application development
    \item Enterprise-level features and integrations
\end{enumerate}

\section{Conclusion}

This review presents a comprehensive analysis of an advanced AI chatbot system that successfully integrates multiple technologies to provide a robust conversational AI platform. The system demonstrates significant capabilities in multi-modal file processing, intelligent conversation management, and specialized tools like resume generation. The proposed enhancements would further strengthen the system's utility and applicability across various domains.

From an academic perspective, the work contributes a unified architecture and design rationale for integrating multimodal conversational AI with ATS-aware resume generation, providing a foundation for comparative evaluation in future studies. The modular architecture and extensible design provide a solid foundation for future development. With deeper empirical evaluation and domain-specific benchmarking, the system can be extended into thesis- or doctoral-level research on AI-assisted career development systems.

\section{References}

\begin{thebibliography}{99}

\bibitem{adamopoulou2020}
Adamopoulou, E., \& Moussiades, L. (2020). Chatbots: History, technology, and applications. \textit{Machine Learning with Applications}, 2, 100006. \url{https://doi.org/10.1016/j.mlwa.2020.100006}

\bibitem{radford2021}
Radford, A., et al. (2021). Learning transferable visual models from natural language supervision. \textit{International Conference on Machine Learning}, 8748-8763.

\bibitem{weizenbaum1966}
Weizenbaum, J. (1966). ELIZA—a computer program for the study of natural language communication between man and machine. \textit{Communications of the ACM}, 9(1), 36-45.

\bibitem{vaswani2017}
Vaswani, A., et al. (2017). Attention is all you need. \textit{Advances in Neural Information Processing Systems}, 30, 5998-6008.

\bibitem{thoppilan2022}
Thoppilan, R., et al. (2022). LaMDA: Language models for dialog applications. \textit{arXiv preprint arXiv:2201.08239}.

\bibitem{baltrusaitis2019}
Baltrusaitis, T., Ahuja, C., \& Morency, L. P. (2019). Multimodal machine learning: A survey and taxonomy. \textit{IEEE Transactions on Pattern Analysis and Machine Intelligence}, 41(2), 423-443.

\bibitem{lu2019}
Lu, J., Batra, D., Parikh, D., \& Lee, S. (2019). ViLBERT: Pretraining task-agnostic visiolinguistic representations for vision-and-language tasks. \textit{Advances in Neural Information Processing Systems}, 32.

\bibitem{devlin2018}
Devlin, J., Chang, M. W., Lee, K., \& Toutanova, K. (2018). BERT: Pre-training of deep bidirectional transformers for language understanding. \textit{arXiv preprint arXiv:1810.04805}.

\bibitem{w3c2023}
W3C. (2023). Web Speech API Specification. \url{https://w3c.github.io/speech-api/}

\bibitem{pradhan2018}
Pradhan, A., Mehta, K., \& Findlater, L. (2018). "Accessibility came by accident": Use of voice-controlled intelligent personal assistants by people with disabilities. \textit{Proceedings of the 2018 CHI Conference on Human Factors in Computing Systems}, 1-13.

\bibitem{black2020}
Black, J. S., \& van Esch, P. (2020). AI-enabled recruiting: What is it and how should a manager use it? \textit{Business Horizons}, 63(2), 215-226.

\bibitem{google2024}
Google AI. (2024). Gemini: A Family of Highly Capable Multimodal Models. \url{https://deepmind.google/technologies/gemini/}

\bibitem{mozilla2024}
Mozilla. (2024). PDF.js: A Portable Document Format (PDF) viewer. \url{https://mozilla.github.io/pdf.js/}

\bibitem{cohere2024}
Cohere AI. (2024). Cohere API Documentation. \url{https://docs.cohere.com/}

\bibitem{resumai2024}
Khan, M., et al. (2024). ResumAI: Revolutionizing Automated Resume Analysis and Recommendation with Multi-Model Intelligence. \textit{IEEE Conference Publication}. \url{https://ieeexplore.ieee.org/document/10426042/}

\bibitem{atsllm2024}
Khadilkar, A. (2024). Optimizing Resume Authenticity and ATS Compatibility with LLM Feedback Integration. \textit{Cal Poly CENG SURP}. \url{https://digitalcommons.calpoly.edu/ceng_surp/72/}

\bibitem{resumebuilder2024}
Sharma, R., et al. (2024). AI-Powered Resume Builder: Enhancing Job Applications with Artificial Intelligence. \textit{IEEE Conference Publication}. \url{https://ieeexplore.ieee.org/document/10986431/}

\end{thebibliography}

\section*{Appendix}

\subsection*{System Requirements}

\begin{itemize}
    \item \textbf{Browser:} Modern browsers (Chrome 90+, Firefox 88+, Safari 14+, Edge 90+)
    \item \textbf{JavaScript:} ES6+ support required
    \item \textbf{Internet Connection:} Required for API calls
    \item \textbf{Storage:} LocalStorage support for chat history
\end{itemize}

\subsection*{API Keys and Configuration}

The system requires the following API keys:
\begin{itemize}
    \item Google Gemini API Key
    \item Cohere API Key (for resume builder features)
\end{itemize}

\subsection*{File Format Support}

\begin{itemize}
    \item \textbf{Images:} JPEG, PNG, GIF, WebP
    \item \textbf{Videos:} MP4, WebM
    \item \textbf{Documents:} PDF, TXT, CSV
    \item \textbf{Maximum File Size:} 10MB (recommended)
\end{itemize}

\vspace{2cm}

\noindent\textbf{Author Information:}
\begin{itemize}
    \item \textbf{Project:} AI Chatbot with Multi-Modal Processing
    \item \textbf{Version:} 1.0
    \item \textbf{Date:} 2024
    \item \textbf{License:} Educational/Research Use
\end{itemize}

\vspace{1cm}

\noindent\emph{This document is original work created for academic and research purposes. All references are properly cited, and the content is designed to be plagiarism-free.}

\end{document}

